\documentclass[10pt]{article}

\usepackage[T1]{fontenc}
\usepackage[utf8]{inputenc}
\usepackage{mathpazo}
\usepackage{microtype}
\usepackage{amsmath,amssymb}
\usepackage{xcolor}
\usepackage[most]{tcolorbox}
\usepackage{enumitem}
\usepackage[hidelinks]{hyperref}
\usepackage[margin=0.50in]{geometry}

\definecolor{ink}{HTML}{1A1A18}
\definecolor{mut}{HTML}{6B6B66}
\definecolor{acc}{HTML}{7A3B2E}
\definecolor{pale}{HTML}{F4F2ED}

\pagestyle{empty}
\setlength{\parindent}{0pt}
\setlength{\parskip}{0.18em}

\newcommand{\zQ}{0_{\mathbb{Q}}}
\newcommand{\zg}{\mathbf{0}_{g}}
\newcommand{\hd}[1]{\par\vspace{0.10em}\noindent
  {\color{acc}\scriptsize\bfseries\MakeUppercase{#1}}\hspace{0.55em}%
  {\color{acc}\leaders\hrule height 0.3pt\hfill\kern0pt}\par\vspace{0.02em}}

\begin{document}

{\Large\bfseries What a Zero Readout Certifies}\par
\vspace{0.15em}
{\color{mut}\normalsize Zero as the failure locus of retained distinction}\par
\vspace{0.35em}
{\color{mut}\footnotesize Yaoharee Lahtee\ \ \textbullet\ \ ORCID 0009-0005-3861-0626\ \ \textbullet\ \ Open Civil Science Initiative\ \ \textbullet\ \ \href{https://doi.org/10.5281/zenodo.21665100}{doi:10.5281/zenodo.21665100}\hfill\itshape (preprint version)}\par
\vspace{-0.2em}
{\color{ink}\rule{\linewidth}{0.8pt}}\par\vspace{-0.15em}\fontsize{9}{10.4}\selectfont

%% ---------------------------------------------------------------
\hd{the setting}

Difference is primitive: to read is to tell two things apart and retain a finite record. A \textbf{declared structure} $g$ is a finite weighted graph stating which pairs this reader is arranged to compare. A field $\Phi$ assigns a rational to each vertex. The retained information of $\Phi$ is
\vspace{-0.5em}\[
I_g(\Phi)\;=\;\sum_{(i,j,w)\in g} w\,(\Phi_i-\Phi_j)^2
\;=\;\Phi^{\top}L_R\,\Phi,
\qquad L_R=D_W-W .
\]\vspace{-0.55em}
Everything is over $\mathbb{Q}$. No limit, no completion, no real line.

%% ---------------------------------------------------------------
\hd{the theorem}

\begin{tcolorbox}[colback=pale,colframe=ink,boxrule=0.9pt,arc=1pt,left=9pt,right=9pt,top=4pt,bottom=4pt]
Let every weight of $g$ be \textbf{strictly positive}. Then for every field $\Phi$:
\[
I_g(\Phi)=\zQ
\iff
\Phi_i=\Phi_j \text{ on every edge}
\iff
\Phi \text{ is constant on every component.}
\]
\end{tcolorbox}

%% ---------------------------------------------------------------
\hd{the proof, complete}

\begin{enumerate}[leftmargin=1.45em,itemsep=0pt,topsep=0pt,parsep=0pt]
\item Each term $w(\Phi_i-\Phi_j)^2$ is a nonnegative weight times a square, hence $\ge 0$.
\item A sum of nonnegatives is $0$ only if every term is $0$. \hfill{\color{mut}\itshape(induction on the edge list)}
\item $w\,d^2=0$ with $w>0$ forces $d^2=0$, and $d^2=0$ forces $d=0$. \hfill{\color{mut}\itshape($\mathbb{Q}$ is an integral domain)}
\item Equality propagates along paths; conversely a constant field kills every term.
\end{enumerate}

{\color{mut}Strict positivity is spent exactly once, at step 3. One edge of weight $0$ breaks the theorem: it contributes nothing while its endpoints differ.}

%% ---------------------------------------------------------------
\hd{two examples that fix the meaning}

\begin{tcolorbox}[colback=white,colframe=mut,boxrule=0.4pt,arc=1pt,left=8pt,right=8pt,top=3.5pt,bottom=3.5pt]
\textbf{Zero on a field that is nowhere zero.}\ \ Let $\Phi_i=17$ everywhere on a connected $g$. Then $I_g(\Phi)=\zQ$. Nothing is empty, no value is zero, nothing is absent --- what is absent is any \emph{difference} for the structure to retain.

\vspace{0.35em}
\textbf{Zero across a gap in the declaration.}\ \ Let $g$ have components $A,B$ with $\Phi=0$ on $A$ and $\Phi=17$ on $B$. Then $I_g(\Phi)=\zQ$ although $\Phi$ is not constant. Where $g$ supplies no path it supplies no comparison; the reader was never arranged to tell those apart.
\end{tcolorbox}

%% ---------------------------------------------------------------
\hd{what the theorem licenses}

Two objects, marked apart, because conflating them is a category mistake:
\vspace{-0.6em}\[
\zQ \in \mathbb{Q}\ \ \text{(a value in the \emph{codomain})}
\qquad\qquad
\zg := I_g^{-1}(\{\zQ\})\ \ \text{(its fibre in the \emph{domain})}
\]\vspace{-0.55em}

\begin{tcolorbox}[colback=white,colframe=acc,boxrule=0.9pt,arc=1pt,left=9pt,right=9pt,top=4pt,bottom=4pt]
\textbf{Claim.} The foundational content of a readout returning $\zQ$ is membership in $\zg$: \emph{the class of fields across which the declared structure retains no difference.} Exact in both directions, by the theorem --- not by definition.
\end{tcolorbox}

\vspace{0.2em}
\textbf{Not claimed.} That $\zQ$ fails to be a mathematical object --- it is a rational, and in a generated tower also the ground of $\mathbb{N}$. That this displaces any account of arithmetic zero --- it answers a different question. That the operator is faithful \emph{absolutely} --- example two shows it is faithful relative to $g$, which is the strongest form available to a reader with a finite declaration.

\textbf{Also not this.} $\zQ$ is a fact about the states: no difference is retained. The unresolved element $\bot$ is a fact about the reader: it cannot tell. Only the first is determinate, which is why only the first admits a biconditional.

%% ---------------------------------------------------------------
\hd{status}

{\scriptsize
The mathematics is elementary and classical over $\mathbb{R}$ (Dirichlet energy vanishes iff constant on components; Chung, Godsil--Royle). \textbf{No novelty is claimed against it.} What is offered is its placement as a \emph{faithfulness condition}: over $\mathbb{Q}$, with the hypothesis explicit, machine-checked without classical axioms, and shown to coincide in both directions with the failure locus of the primitive --- closing a join that was previously carried by interpretation.

\vspace{0.3em}
Verified in Coq: 17 named results, \verb|Print Assumptions| $=$ \emph{Closed under the global context} on each. No \verb|Reals|, no \verb|classic|. Source: \href{https://github.com/morrocwi/zero-readout-certifies}{github.com/morrocwi/zero-readout-certifies}\ \ \textbullet\ \ archived at \href{https://doi.org/10.5281/zenodo.21665100}{doi:10.5281/zenodo.21665100}. Don't take the proof on trust --- compile it.
}

\vspace{0.15em}
{\color{ink}\rule{\linewidth}{0.4pt}}
\vspace{0.05em}
{\color{mut}\scriptsize\itshape Arithmetic supplies the symbol $\zQ$; the retained-difference theorem determines what a legitimate occurrence of that symbol certifies.}

\end{document}
